\scanpage{67}
\mspage{5.}
% archive sheet number: 66

3.2. \uline{Les missions d'enseignement et les missions de recherche}

- Les établissements ont-ils des missions de recherche ? C'est là une question à nous poser et à laquelle il importe de répondre en allant au fond du problème.

En effet, la recherche est aujourd'hui une activité fondamentale pour la nation ; des milliers de travailleurs y consacrent la totalité de leur activité, dans le cadre du"plan recherche". On pourrait se demander s'il est nécessaire que les enseignants des établissements d'enseignement supérieur y consacrent aussi une part de leur activité... Il n'est pas évident que l'on puisse faire du bon travail à temps partiel à l'époque d'intense compétition internationale qu'est la notre. Vaudrait-il mieux peut être, développer un peu plus l'effort de recherche et découpler totalement l'enseignement et la recherche, à l'époque où l'enseignement supérieur ne peut plus revendiquer un monopole dans cette activité ?...

On ne peut ignorer que la question se pose en ces termes à de nombreuses personnes. \unsure{-}

- Face à cette intérrogation, la D.E.S.U.R. a donc défini très clairement des objectifs spécifiques pour la recherche menée dans les établissements qui doivent donc s'appeler établissements d'enseignement supérieur et de recherche. La base du raisonnement est la suivante :

\noindent\fbox{\parbox{0.95\linewidth}{Une centaine d'établissements autonomes, dont la plupart ont leur siège social en province, exercent en priorité des missions d'enseignement. Toutefois, leur personnel enseignant, formé par la recherche, est disponible à temps partiel pour contribuer à l'effort national de recherche. Du fait même qu'il est enseignant, il peut remplir des missions de recherche spécifiques, indispensabl\unsure{\={e}|e}s à l'effort national de recherche. Cet effort de recherche permet sa formation permanente et la formation de jeunes chercheurs\handwritten{,} contribu\deleted{\unsure{e|a}}ant ainsi à maintenir l'établissement en pleine activité créatrice.}}

Il faut définir, à partir de ces données, des objectifs de recherche pour les établissements d'enseignement su\deleted{\unsure{ép|pé}}périeur. Ceux-ci pourront tout en contribuant à l'effort national de recherche, se développer eux-mêmes pour le plus grand bien de leur mission fondamentale d'enseignement.
